\documentclass[a4paper,10pt]{article}
\usepackage[czech]{babel}
\usepackage[T1]{fontenc}
\usepackage[utf8]{inputenc}
\usepackage{amsmath,amssymb}
\usepackage[margin=18mm]{geometry}
\setlength{\parindent}{0pt}
\begin{document}
\section*{Geometrická posloupnost E/F -- řešení a bodování}
Podklad pro učitele k testu
\texttt{01\_test\_GeometrickaPosloupnost\_EF.tex}.
Zadání jsou doslova převzata z pracovního listu
\texttt{4\_rocnik/pracovni\_listy/pracovniListGeoPosl.tex}.

Zařazení: 4. ročník SPŠ; tematický plán
\texttt{4\_rocnik/tematicke\_plany/SPS.pdf}, s.~1, září, bod 1.3.
Celek posloupností má v září společnou dotaci 12 hodin.
Test používá obecný člen, vztah mezi dvěma členy, konečný součet
a vkládání kladných členů. Nevyžaduje další témata ani kalkulačku.
Odhad doby řešení včetně kontroly je 25--30 minut; není ověřen ve třídě.

\begin{center}
\begin{tabular}{c|c|c|l}
Úloha testu & Varianta E & Varianta F & Ověřovaná dovednost\\ \hline
1 & 6 & 7 & Rozpoznání geometrické posloupnosti z obecného členu\\
2 & 21 & 22 & Určení kvocientu a dalších členů ze dvou členů\\
3 & 40 & 39 & Konečný součet se záporným kvocientem\\
4 & 61 & 62 & Vložení kladných členů mezi dvě zadaná čísla
\end{tabular}
\end{center}

\subsection*{Varianta E}
\begin{enumerate}
\item \textbf{Původní úloha 6.}
Pro každé $n\in\mathbb N$, $n\geq1$, je $a_n=(-2)^{n+1}\ne0$.
Můžeme tedy vytvořit podíl
\[
\frac{a_{n+1}}{a_n}=\frac{(-2)^{n+2}}{(-2)^{n+1}}=-2.
\]
Podíl je konstantní pro všechna přípustná $n$, posloupnost je geometrická a
\[
\boxed{a_1=4,\qquad q=-2.}
\]
Nezávislé ověření: $(-2)^{n+1}=4(-2)^{n-1}=a_1q^{n-1}$.

\item \textbf{Původní úloha 21.}
Nenulové členy $a_3=12$, $a_6=96$ vylučují $a_1=0$ i $q=0$.
Vztah $a_6=a_3q^3$ dává
\[
96=12q^3\iff q^3=8\iff q=2.
\]
Rovnice $q^3=8$ má právě jeden reálný kořen;
hodnota $q=-2$ podmínku nesplňuje. Dále
\[
a_1=\frac{a_3}{q^2}=\frac{12}{4}=3,\qquad
a_9=a_6q^3=96\cdot8=768.
\]
Jediné řešení je $\boxed{q=2,\ a_1=3,\ a_9=768}$.
Zkouška postupným násobením:
\[
3,6,12,24,48,96,192,384,768.
\]
Třetí, šestý a devátý člen souhlasí s použitými hodnotami.

\item \textbf{Původní úloha 40.}
Je dáno $a_1=6$, $q=-\frac12$. Protože $q\ne1$, platí vzorec
$S_n=a_1\frac{1-q^n}{1-q}$. Dosadíme
\[
S_{10}=6\frac{1-\left(-\frac12\right)^{10}}{1-\left(-\frac12\right)}
=6\frac{1-\frac1{1024}}{\frac32}
=4\cdot\frac{1023}{1024}
=\boxed{\frac{1023}{256}}.
\]
Nezávislá zkouška přímým součtem deseti členů:
\[
6-3+\frac32-\frac34+\frac38-\frac3{16}
+\frac3{32}-\frac3{64}+\frac3{128}-\frac3{256}
=\frac{1536-768+384-192+96-48+24-12+6-3}{256}
=\frac{1023}{256}.
\]

\item \textbf{Původní úloha 61.}
Všech pět čísel je kladných, proto $q>0$. První člen je $a_1=2$
a pátý $a_5=162$, tedy
\[
162=2q^4\iff q^4=81.
\]
Z reálných kořenů $q=3$ a $q=-3$ je přípustný pouze $q=3$;
záporný kvocient by vedl k zápornému druhému a čtvrtému členu.
Vkládáme tedy
\[
\boxed{6,18,54},\qquad\text{celá pětice je }2,6,18,54,162.
\]
Zkouška: $\frac62=\frac{18}6=\frac{54}{18}=\frac{162}{54}=3$;
vložené členy jsou kladné. Kladný kořen rovnice pro $q$ je jediný,
takže další přípustné řešení neexistuje.
\end{enumerate}

\newpage
\subsection*{Varianta F}
\begin{enumerate}
\item \textbf{Původní úloha 7.}
Pro každé $n\in\mathbb N$, $n\geq1$, je $2^n>0$ a $a_n=\frac5{2^n}\ne0$.
Platí
\[
\frac{a_{n+1}}{a_n}=\frac{5}{2^{n+1}}\frac{2^n}{5}=\frac12.
\]
Podíl je konstantní pro všechna přípustná $n$, posloupnost je geometrická a
\[
\boxed{a_1=\frac52,\qquad q=\frac12.}
\]
Nezávislé ověření:
$\frac5{2^n}=\frac52\left(\frac12\right)^{n-1}=a_1q^{n-1}$.

\item \textbf{Původní úloha 22.}
Z $a_2=-6\ne0$, $a_5=48\ne0$ plyne $a_1\ne0$, $q\ne0$.
Vztah $a_5=a_2q^3$ dává
\[
48=-6q^3\iff q^3=-8\iff q=-2.
\]
Jde o jediný reálný kvocient. Pak
\[
a_1=\frac{a_2}{q}=\frac{-6}{-2}=3,\qquad
a_8=a_5q^3=48(-8)=-384.
\]
Jediné řešení je $\boxed{q=-2,\ a_1=3,\ a_8=-384}$.
Zkouška postupným násobením:
\[
3,-6,12,-24,48,-96,192,-384.
\]
Druhý, pátý a osmý člen souhlasí s použitými hodnotami.

\item \textbf{Původní úloha 39.}
Je dáno $a_1=3$, $q=-2\ne1$, proto
\[
S_9=3\frac{1-(-2)^9}{1-(-2)}
=3\frac{1+512}{3}=\boxed{513}.
\]
Nezávislá zkouška přímým součtem devíti členů:
\[
3-6+12-24+48-96+192-384+768=513.
\]
Na rozdíl od sudé mocniny v úloze E3 je zde $(-2)^9$ záporné.

\item \textbf{Původní úloha 62.}
Všech šest čísel je kladných, proto $q>0$. Máme $a_1=3$, $a_6=96$ a
\[
96=3q^5\iff q^5=32\iff q=2.
\]
Rovnice $q^5=32$ má právě jeden reálný kořen, který je kladný.
Vkládáme tedy
\[
\boxed{6,12,24,48},\qquad\text{celá šestice je }3,6,12,24,48,96.
\]
Zkouška:
$\frac63=\frac{12}6=\frac{24}{12}=\frac{48}{24}=\frac{96}{48}=2$.
Všechny vložené členy jsou kladné a jiné přípustné řešení neexistuje.
\end{enumerate}

\newpage
\subsection*{Orientační bodování pro obě varianty}
Nezávazný návrh pouze pro učitele. Uznávejte i jiné správné postupy;
bodování a kritéria lze upravit podle odborného posouzení.
Stejnou početní chybu netrestejte opakovaně, pokud na ni navazuje
jinak správný postup.
\begin{enumerate}
\item 3 body: obecný podíl nebo ekvivalentní vzorec pro geometrickou
posloupnost (1), první člen (1), kvocient a závěr pro všechna přípustná $n$ (1).
\item 3 body: kvocient ze vztahu mezi známými členy (1), první člen (1),
hodnota požadovaného členu ($a_9$ ve variantě E, $a_8$ ve variantě F) (1).
\item 3 body: vzorec pro součet geometrické řady (1), správné dosazení
a úprava mocnin včetně znamének (1), výsledek (1).
\item 4 body: vztah mezi krajními členy se správným exponentem (1),
kladný kvocient (1), správný počet vložených členů (1),
správné hodnoty všech vložených členů (1).
\end{enumerate}
Celkem 13 bodů za každou variantu. Zkoušky uvedené v řešení slouží
ke kontrole; žákovské zadání je výslovně nepožaduje, proto za jejich
neprovedení neodečítejte body. Orientační doby řešení úloh jsou
3, 6, 7 a 5 minut; s kontrolou přibližně 25--30 minut.

\subsection*{Srovnatelnost a kontrola}
První dvojice vyžaduje podíl obecných členů a práci s mocninami.
Ve druhé dvojici vede rozdíl indexů vždy na třetí mocninu kvocientu
a jediné reálné řešení; F2 navíc pracuje se záporným kvocientem.
Ve třetí dvojici se používá stejný vzorec se záporným kvocientem.
E3 je početně o něco náročnější kvůli zlomkům, F3 pracuje s celými čísly.
Ve čtvrté dvojici se hledá jediný kladný kořen rovnice pro kvocient;
E4 vyžaduje odmítnutí záporného kořene, F4 o jeden vložený člen více.
Rozdíly jsou malé a obě varianty mají přibližně stejný rozsah i časovou
náročnost. Výběr respektuje zadané úlohy běžného testu;
obtížnost zde není navyšována nad ně pouze kvůli označení E/F.

Každé z osmi zadání bylo ověřeno samostatně. U obecných členů
je podíl doplněn totožností s tvarem $a_1q^{n-1}$; známé členy,
součty a vložená čísla jsou kontrolovány přímým generováním posloupností,
přesnou racionální aritmetikou a dosazením do původních podmínek.
Úplnost řešení plyne z jednoznačnosti reálných lichých odmocnin
a kladného kořene ve čtvrté mocnině. Před každým dělením je
zajištěna nenulovost příslušného výrazu; při použití vzorce pro součet
je splněno $q\ne1$.
\end{document}
